\begin{proof}[Proof of \cref{obj:lem:same-cell-factorial-cross-moment}]
Write
\[
p:=p_k=P(X=k).
\]
For an observed record \(O=(X,A,Y)\), write \(V=Y/M\).  Define the left coordinate factors
\[
L_q(O):=
\begin{cases}
V\mathbf 1\{X=k,A=a\}, & q=0,\\
\mathbf 1\{X=k\}, & q=1,\\
\mathbf 1\{X=k,A=a\}, & 2\le q\le j+1,
\end{cases}
\]
and define the right coordinate factors \(R_s\) in the same way, with \((b,r,s)\) in place of \((a,j,q)\).

For a partial matching \(N\) of size \(h\) between \(\{0,\ldots,j+1\}\) and \(\{0,\ldots,r+1\}\), let \(T_N\) be the merged coordinate set, and let \(\iota_L\) and \(\iota_R\) map left and right coordinates into \(T_N\). Thus \(\iota_R(s)=\iota_L(q)\) exactly when \(N\) matches \(q\) with \(s\). Put
\[
\gamma_N:=
\frac{(m)_{(j+2)+(r+2)-h}}{(m)_{j+2}(m)_{r+2}},
\]
\[
F_{N,t}(O):=
\prod_{q:\,\iota_L(q)=t}L_q(O)
\prod_{s:\,\iota_R(s)=t}R_s(O),
\qquad
M_N:=\int \prod_{t\in T_N}F_{N,t}(O_t)\,
  dP_O^{\otimes T_N}((O_t)_{t\in T_N}).
\]
Finally set
\[
m_{k,a,j}:=\int\prod_{q=0}^{j+1}L_q(O_q)\,
  dP_O^{\otimes(j+2)}((O_q)_q),
\qquad
m_{k,b,r}:=\int\prod_{s=0}^{r+1}R_s(O_s)\,
  dP_O^{\otimes(r+2)}((O_s)_s).
\]
The block-size condition implies \(K\le m\), hence \(j+2\le m\) and \(r+2\le m\). The selector factors are measurable bounded indicators, and the marked factor is integrable under the moment conditions in \cref{obj:def:model-class}; the merged kernels are finite products with at most two marked factors at any single observed coordinate, hence integrable. With the coordinate factors just defined, each statistic \(U_{k,a,j}\) and \(U_{k,b,r}\) is its normalized ordered-product sum over injective maps into the block. Expanding the product of the two normalized sums and grouping each ordered pair of injective maps by the partial matching \(N\) that records their shared block indices gives a contribution \(\gamma_N M_N\) for matching pattern \(N\). The empty matching has product moment \(m_{k,a,j}m_{k,b,r}\) by independence and is then centered by subtracting \(m_{k,a,j}m_{k,b,r}\). Therefore
\[
\operatorname{Cov}(U_{k,a,j},U_{k,b,r})
=
(\gamma_\varnothing-1)m_{k,a,j}m_{k,b,r}
+
\sum_{h=1}^{\min(j+2,r+2)}
\sum_{N\in\mathcal M_{j,r}(h)}
\gamma_N\,M_N,
\]
where \(\mathcal M_{j,r}(h)\) is the finite set of partial matchings of size \(h\).
% lean: centeredCrossMoment_allBlockOrderedMarkedFactorial

Each one-coordinate mean is bounded by the cell mass:
\[
\left|\int L_q(O)\,dP_O(O)\right|\le p,
\qquad
\left|\int R_s(O)\,dP_O(O)\right|\le p.
\]
For the outcome-marked coordinate, the argument splits according to the cell mass. If \(p=0\), the observed arm-cell event \(\{X=k,A=c\}\) has \(P_O\)-mass zero for each arm \(c\), so the marked integral is zero. If \(p>0\), the supported-cell branch of \cref{obj:ass:mean-normalization} applies to the arm-cell outcome law for each arm \(c\), and \cref{obj:ass:overlap} supplies positive observed arm-cell mass when this law is read as the observed conditional law. The mean bound and \(M\ge1\) from \cref{obj:def:model-class} give absolute normalized arm-cell mean at most \(1/2\). The observed arm-cell restriction has total mass
\(p\{\pi_k\mathbf 1_{\{c=1\}}+(1-\pi_k)\mathbf 1_{\{c=0\}}\}\le p\), so the absolute marked integral is bounded by \(p/2\). The selector-only coordinates integrate to cell or arm-cell probabilities bounded by \(p\). Since the product-law means factor coordinatewise,
\[
|m_{k,a,j}|\le p^{j+2},
\qquad
|m_{k,b,r}|\le p^{r+2}.
\]
The empty-matching normalization satisfies, under \(j+2\le K\), \(r+2\le K\), and \(4(K+2)^2\le m\),
\[
|\gamma_\varnothing-1|
\le
\frac{2(K+2)^2}{m}.
\]
Therefore
\[
|(\gamma_\varnothing-1)m_{k,a,j}m_{k,b,r}|
\le
\frac{2(K+2)^2}{m}p^{j+2}p^{r+2}.
\]
% lean: marked_emptyMatchingNormalization_sub_one_le
% lean: orderedProductMean_markedFactorialCoordinate_abs_le

It remains to bound the positive-matching part.  For each merged coordinate \(t\in T_N\), the fiber definition of \(F_{N,t}\) gives exactly one of the following forms: a single left factor \(L_q\), a single right factor \(R_s\), or a matched product \(L_qR_s\). The first two cases include the lone unmarked cell-indicator coordinate. For single factors, the preceding one-coordinate bound gives
\[
\left|\int L_q\,dP_O\right|\le p\le \frac54p,
\qquad
\left|\int R_s\,dP_O\right|\le p\le \frac54p.
\]
For matched products, first note the square-moment bounds
\[
\int L_q(O)^2\,dP_O(O)\le \frac54p,
\qquad
\int R_s(O)^2\,dP_O(O)\le \frac54p.
\]
For the outcome-marked coordinate these bounds again split on \(p\): the zero-mass branch is null on every arm-cell event, while on \(p>0\), \cref{obj:ass:overlap} supplies the positive arm-cell support for either arm. The arm-cell law then satisfies the mean and central second-moment bounds from \cref{obj:ass:mean-normalization,obj:ass:second-central-moment}; expanding the second moment around \(\mu_{ck}\) gives an unnormalized second moment at most \(5M^2/4\), hence, using \(M\ge1\) from \cref{obj:def:model-class}, a normalized second moment at most \(5/4\). Multiplying by the observed arm-cell mass, which is at most \(p\), gives the displayed square-moment bound for the marked coordinate. Selector coordinates are idempotent indicators with event mass at most \(p\). Hence, by integrability, \(|\int fg|\le\int|fg|\) and \(2|fg|\le f^2+g^2\),
\[
\left|\int L_q(O)R_s(O)\,dP_O(O)\right|
\le
\frac12\int\{L_q(O)^2+R_s(O)^2\}\,dP_O(O)
\le
\frac54p.
\]
Consequently every merged coordinate satisfies
\[
\left|
\int F_{N,t}(O)\,dP_O(O)
\right|
\le
\frac54p.
\]
% lean: mergedMarkedCoordinateFactor_integral_abs_le

The merged product moment factors over the merged coordinates:
\[
M_N=\prod_{t\in T_N}\int F_{N,t}(O)\,dP_O(O).
\]
Since \(|T_N|=(j+2)+(r+2)-h\),
\[
|M_N|
\le
\left(\frac54p\right)^{(j+2)+(r+2)-h}.
\]
% lean: mergedProductMoment_markedFactorialCoordinate_same_cell_abs_le

Using \(\gamma_N\ge0\) and the triangle inequality,
\[
\left|
\sum_{h=1}^{\min(j+2,r+2)}
\sum_{N\in\mathcal M_{j,r}(h)}
\gamma_NM_N
\right|
\le
\sum_{h=1}^{\min(j+2,r+2)}
\sum_{N\in\mathcal M_{j,r}(h)}
\gamma_N
\left(\frac54p\right)^{(j+2)+(r+2)-h}.
\]
It remains only to sum the normalization weights.  Let \(u=j+2\) and \(v=r+2\). For each \(h\), the size-\(h\) normalization bound and the partial-matching count give
\[
\sum_{N\in\mathcal M_{j,r}(h)}\gamma_N
\le
\binom{u}{h}v^h\frac{\exp(1)}{m^h}.
\]
Here \(|\mathcal M_{j,r}(h)|=\binom{u}{h}\binom{v}{h}h!\), and the displayed bound uses \(\binom{v}{h}\le v^h/h!\) together with the factorial-normalization estimate supplied by the block-size condition \(4(K+2)^2\le m\). Therefore, for every \(x\ge0\),
\[
\begin{aligned}
&\sum_{h=1}^{\min(u,v)}
\sum_{N\in\mathcal M_{j,r}(h)}
\gamma_N x^{u+v-h} \\
&\qquad\le
\sum_{h=1}^{\min(u,v)}
\exp(1)x^v\binom{u}{h}x^{u-h}\left(\frac{v}{m}\right)^h \\
&\qquad\le
\sum_{h=0}^{u}
\exp(1)x^v\binom{u}{h}x^{u-h}\left(\frac{v}{m}\right)^h
=
\exp(1)x^v\left(x+\frac{v}{m}\right)^u.
\end{aligned}
\]
The choice \(x=(5/4)p\) is admissible because \(p\ge0\), and it yields
\[
\left|
\sum_{h=1}^{\min(j+2,r+2)}
\sum_{N\in\mathcal M_{j,r}(h)}
\gamma_NM_N
\right|
\le
\exp(1)\left(\frac54p\right)^{r+2}
\left(\frac54p+\frac{r+2}{m}\right)^{j+2}.
\]
% lean: positiveMatchingNormalization_weighted_sum_le

Combining the empty-matching and positive-matching bounds by the triangle inequality gives
\[
\left|\operatorname{Cov}_{Q_P^{(m)}}(U_{k,a,j},U_{k,b,r})\right|
\le
\frac{2(K+2)^2}{m}p_k^{j+2}p_k^{r+2}
+
\exp(1)\left(\frac54p_k\right)^{r+2}
\left(\frac54p_k+\frac{r+2}{m}\right)^{j+2},
\]
as claimed.
% lean: centeredCrossMoment_allBlockOrderedMarkedFactorial_same_cell_abs_le
\end{proof}
